% Nouveau EGA V (cote 25): archive cote 25, site page 163 (PDF page 164).
% typescript. Draft compared: none.
% AI-transcribed by Claude Opus 5.5 (claude-opus-5-5, adaptive thinking, effort high), Message Batches API, batch ega25b (msgbatch_01A7n5YuXubUPNK1zsepi8RG), 2026-11-13; not yet checked by eye.
\scanpage{25/163}
\mspage{\hand{6}}

un sous-préschéma X" de X, régulier,  contenant X, et tel que X
puisse se décrire dans X" par une seule équation $\phi = 0$ , admettant x comme
zéro sing élémentaire de multiplicité n.

On conclut en particulier que si A est local noeth élémentaire de

mult. n, alors A est Coh. Mac.  Nous montrerons au par.20

(par un calcul facile d'éclatement) que le point fermé de Spec(A) est
le seul point singulier de Spec(A); il en résulte que  A est normal
sss dim A $\geqslant$ 2, réduit sss dimA $\geqslant$1. Si dim A = 0, alors  A est élémentaire
de multiplicité n sss  $m/m^2$ est de \unsure{rang}1 et  n est le plus petit entier
tel que $\underline{m}^n = 0$; on notera que ces anneaux sont aussi lesquotients d'anneaux
de valuion discrète par la puissance n.ème de leur idéal maximal.

15.7.   Il y a lieu de dégager les variantes ``géométriques'' de 15.6. ,
on trouvera en particulier que si X est un préschéma loc de type fini sur
le corps k, $x \in X$, et si d $= \dim_x X$ , alors  x est une singularité 
élémentaire
 de multiplicité  n de X, sss il existe un voisinage ouvert U de x
tel que U se plonge comme sous-préschéma d'un préschéma lisse X' sur k,
 conne xe et de dimension d+1 , défini par une équation $\phi$=0 , x étant
un zéro géom. singulier ordinaire de multiplicité n de $\phi$ . On aura dit d'ailleurs
\placedmargin{$\surd$} dans 15.5. que cela signifie  que  la valeur en x de partie principale $d^n\phi$ de $\phi$ , 
 (qui à priori est un\unsure{e}  élément de $\underline{P}^n_{X'/k}\otimes \unsure{\underline{k}(x)}$ ,et plus
précisément de son Idéal d'augmentation $\underline{I}$ )  est en fait un élément de
 $\underline{I}^n = \underline{Sym}^n_{\underline{O}_U}(\underline{\Omega}^1_{X'/k}\otimes k(\hand{x}))$), et  à ce titre est une  forme lisse 
\placedmargin{|} \hand{[}NB Oubli dans 15.1. : l'ensemble des points de lissité d'une forme est
\placedmargin{|} ouvert\hand{]}. On remarquera de même qu'un tel point est isolé dans l'ensemble
des points non lisses, il est géom normal  (resp. géom réduit) sss
d $\geqslant 2$  (resp. d $\geqslant 1$). Si d = 0 et k(x) = k (i.e. x  isolé et x rationnel sur k) alors 
point singulier élémentaire  de mult\hand{,} \unsure{n} signifie que $\underline{O}X,x$  est  k-

\mspage{\hand{163}}
